% =====================================================================
%  Chapter 30 station file -- 14 chaptersections + Summit + lastword.
%    1 Summary                8 Recommended Study Order
%    2 Master Decision Flow   9 Difficulty Progression
%    3 Common Mistakes       10 Reduction Gallery
%    4 WA Clinic             11 Mixed Problem Set (10 practice items)
%    5 Quick Reference       12 Summit -- CF 917D Stranger Trees
%    6 Write-From-Memory     13 Tool Kit Inventory
%    7 Complete Checklist    14 Related Topics
%  Summit live-verified: title "D. Stranger Trees", 1s TL, 256 MB, n<=100,
%  mod 10^9+7.  Samples: star on 3 -> "0 2 1"; path on 4 -> "1 7 7 1";
%  star on 4 -> "0 9 6 1".  All three hand-replayed in the Trace section.
% =====================================================================

\chaptersection{Chapter Summary}

Chapter~30 is the enumerative counterweight to Chapters 12--29.  Where those
chapters \emph{operate} on a tree this one \emph{counts} trees.  Four tools
carry the chapter:

\begin{itemize}[nosep]
  \item the \textbf{Pr\"ufer bijection} (\(\secref{c32-prufer}\)):
        \(\mathcal{T}_n\leftrightarrow[n]^{n-2}\), with the code's multiplicity
        of label \(v\) equal to \(\deg_T(v)-1\);
  \item \textbf{Cayley's formula and its variants} (\(\secref{c32-cayley}\)):
        \(n^{n-2}\), rooted \(n^{n-1}\), degree multinomial \((n-2)!/\prod(d_i-1)!\),
        forest-completion \(n^{k-2}\prod_i s_i\);
  \item the \textbf{Matrix-Tree theorem} (\(\secref{c32-matrix}\)): spanning
        trees of a given graph \(G\) equal any cofactor of \(L(G)=D-A\); a
        modular determinant is the algorithm;
  \item \textbf{linearity of expectation} (\(\secref{c32-expect}\)): per-vertex
        indicator decomposition, with the CF 280C process collapsing to
        \(\sum_v 1/\mathrm{depth}(v)\).
\end{itemize}

\(\secref{c32-boundary}\) closes the loop by drawing the line between
\emph{labelled} counts (this chapter) and \emph{unlabelled} shapes
(Ch.~28 \chref{c22}), via the orbit-stabiliser bridge
\(\thmref{c32-orbit}\).  The Summit \(\probref{c32-summit}\) is a genuine
labelled-count question mark whose answer uses the forest-completion formula
as its workhorse and a binomial inversion to close the "exactly \(k\) shared
edges" filter.

\chaptersection{Master Decision Flowchart}

\begin{center}
\begin{tikzpicture}[diagram, scale=0.85, transform shape,
   q/.style={rectangle, rounded corners=3pt, draw=primaryblue!70, fill=primaryblue!8,
             line width=0.8pt, minimum width=26mm, minimum height=8mm,
             font=\scriptsize\bfseries, inner sep=2pt, align=center},
   a/.style={rectangle, rounded corners=2pt, draw=successgreen!70, fill=successgreen!10,
             line width=0.8pt, minimum width=32mm, minimum height=8mm,
             font=\scriptsize, inner sep=2pt, align=center},
   every edge/.style={-Latex, line width=0.8pt, draw=darkgray}]
  \node[q] (root) at (0,0) {Read statement};
  \node[q] (lab)  at (-4,-1.6) {labelled?};
  \node[q] (un)   at (4,-1.6) {given graph\\ \(G\)?};
  \node[a] (pf)  at (-6,-3.4) {Pr\"ufer\\ \(\secref{c32-prufer}\)};
  \node[a] (cy)  at (-2.6,-3.4) {Cayley /\\ completion};
  \node[a] (mt)  at (1.6,-3.4) {Matrix-\\ Tree};
  \node[a] (ex)  at (5.6,-3.4) {Linearity\\ of \(\mathbb{E}\)};
  \node[a] (bo)  at (8,-3.4) {Hand off\\ Ch.28/26/31};
  \draw (root)--(lab)--(pf);
  \draw (lab)--(cy);
  \draw (root)--(un)--(mt);
  \draw (un)--(ex);
  \draw (un.east) to[out=0,in=-90] (bo.south);
\end{tikzpicture}
\end{center}

Every Ch.~30 problem walks one of the four leaf boxes.  The "Hand off"
exit is \(\secref{c32-boundary}\)'s pick list: an unlabelled question marks
a Ch.~28 or Ch.~26 problem dressed as a Ch.~30 one.

\chaptersection{Common Mistakes Checklist}

\begin{enumerate}[nosep]
  \item \textbf{Using Cayley for a \emph{given} graph.}  Cayley \(n^{n-2}\)
        applies only when the graph is \(K_n\); for an arbitrary \(G\) you
        must use \(\thmref{c32-matrix}\) or lose correctness on non-complete
        instances.
  \item \textbf{Dividing a labelled count by \(n!\) to get the unlabelled
        one.}  The orbit sizes are \(n!/|\mathrm{aut}(S)|\), not uniform;
        \(\thmref{c32-orbit}\) has to be summed shape-by-shape.
  \item \textbf{Pr\"ufer code length off by one.}  \(n^{n-2}\) is derived
        from a code of length \(\boldsymbol{n-2}\) over \([n]\); a length
        \(n-1\) code enumerates \emph{rooted} trees.
  \item \textbf{Reading \(P(T)=c_{n-2}\).}  The last-code-entry rule is
        wrong whenever \(c_{n-2}=n\) (see the counterexample in
        \(\secref{c32-cayley}\)).
  \item \textbf{Modular determinant without a prime modulus.}  Gaussian
        elimination requires an invertible pivot; use a prime (\(10^9+7\),
        \(998244353\)) or switch to Bareiss' exact-division variant.
  \item \textbf{Linearity with independence.}  \(\lemref{c32-indexp}\)
        does not assume independence, so a proof that "justifies"
        independence is at best wasted work and at worst hides a bug in
        the decomposition.
  \item \textbf{Depth \(0\) vs.\ depth \(1\) for the root.}  CF 280C's
        expectation is \(\sum_v 1/\mathrm{depth}(v)\) with
        \(\mathrm{depth}(1)=1\); a \(0\)-based DFS depth diverges on the
        root term.
  \item \textbf{Handling \(k=1\) in \(n^{k-2}\).}  The exponent
        becomes \(-1\) at \(k=1\); special-case the "already a tree"
        branch (total \(=1\)).
  \item \textbf{Forgetting the \(s_i\) product in forest-completion.}
        \(n^{k-2}\) alone is the count when each component is a single
        vertex; the product \(\prod s_i\) is the local graft choice.
  \item \textbf{Skipping the binomial inversion in the Summit.}  The DP
        returns \(f(k)\) -- the count of trees \emph{containing} a fixed
        \(k\)-set -- not \(g(k)\), the exact-\(k\) count.  Forgetting the
        inversion produces answers that are systematically too big.
\end{enumerate}

\chaptersection{Wrong-Answer Clinic}

\begin{itemize}[nosep]
  \item \textbf{WA on CF 2190D sample case 1 (n=3, m=0).}  A common cause:
        counting all three labelled trees under the assumption
        \(P(T)=c_{n-2}\).  Codes \([1],[2],[3]\) give \(P(T)=1,2,2\), and
        the "last entry" heuristic assigns \(1,2,3\) -- wrong on the
        third.  The histogram is \((1,2,0)\) instead of the correct
        \((1,2)\) once you exclude \(v=n\).
  \item \textbf{WA on CF 280C sample \(n=2\).}  Emitting \(0.5\) instead
        of \(1.5\): the root \(1\) is at depth \(0\) in a BFS-level view,
        so \(1/\mathrm{depth}\) is either undefined or you have to skip
        the root.  Both are wrong -- use vertex-count depth starting at
        \(1\).
  \item \textbf{WA on the Summit for \(k=0\).}  Reporting \(f(0)=n^{n-2}\)
        as the answer; but \(g(0)\) is not \(f(0)\) -- it is
        \(\sum_j(-1)^j\binom{j}{0}f(j)=\sum_j(-1)^j f(j)\), the
        inclusion-exclusion over all tree edges.  On the star \(n=3\):
        \(f(0)-f(1)+f(2)=3-4+1=0\), matching sample.
  \item \textbf{WA on Matrix-Tree for \(G\) disconnected.}  Your code
        prints the determinant of \(L^{(i)}\) as if it were nonzero; but
        the kernel of a disconnected \(G\)'s Laplacian has dimension \(>1\)
        so every cofactor is zero.  The right answer is \(0\), and any
        nonzero output means you built \(L\) wrong (typically a missing
        degree on the diagonal).
\end{itemize}

\chaptersection{Quick Reference Card}

\begin{quickref}
\begin{itemize}[nosep]
  \item \textbf{Pr\"ufer code.}  Length \(n-2\), labels in \([n]\),
        \(m_v(T)=\deg_T(v)-1\) copies of \(v\).
  \item \textbf{Cayley.}  \(|\mathcal{T}_n|=n^{n-2}\).
  \item \textbf{Rooted.}  \(n^{n-1}\).
  \item \textbf{Degree prescribed.}  \((n-2)!/\prod_i(d_i-1)!\).
  \item \textbf{Forest completion.}  \(n^{k-2}\prod_i s_i\).
  \item \textbf{Complete bipartite \(K_{a,b}\).}  \(a^{b-1}b^{a-1}\).
  \item \textbf{Cycle \(C_n\).}  \(n\) spanning trees.
  \item \textbf{Matrix-Tree.}  \(\tau(G)=\det L^{(i)}(G)\), any \(i\).
  \item \textbf{Linearity.}  \(\mathbb{E}[\sum X_i]=\sum \mathbb{E}[X_i]\),
        no independence needed.
  \item \textbf{CF 280C.}  \(\mathbb{E}=\sum_v 1/\mathrm{depth}(v)\),
        depth in vertices from root.
  \item \textbf{Summit (CF 917D).}  \(f(k)=n^{n-k-2}\sum_S\prod_i s_i\),
        then \(g(k)=\sum_{j\ge k}(-1)^{j-k}\binom{j}{k}f(j)\).
\end{itemize}
\end{quickref}

\chaptersection{Write-From-Memory Drills}

Each drill is a small program you should be able to reproduce on a blank
editor in ten minutes; the labels point back to the code the drill rebuilds.

\begin{enumerate}[nosep]
  \item Encode \([5,5,3,4]\)-style Pr\"ufer code from an edge list --
        rebuild \(\coderef{c32-encode}\) in \(\bigO{n\log n}\) with a
        min-heap of leaves and a degree array.
  \item Decode a Pr\"ufer code back to an edge list -- rebuild
        \(\coderef{c32-decode}\) with a count array and a min-heap of
        zero-count labels; the final two heap entries are \(n\) and \(P(T)\).
  \item Given a Laplacian as a \(2\)D vector, compute the modular
        determinant via Gaussian elimination -- rebuild
        \(\coderef{c32-moddet}\); handle row swaps and pivots explicitly.
  \item Compute \(\sum_v 1/\mathrm{depth}(v)\) on a rooted tree given
        as an edge list -- rebuild \(\coderef{c32-game}\); iterative DFS,
        double accumulator, print with \(20\) digits.
  \item Compute \(n^{k-2}\prod_i s_i\) modulo a prime for a forest --
        rebuild \(\coderef{c32-modpow}\) and its DSU driver from
        \(\coderef{c32-prufvtx}\); special-case \(k=1\).
\end{enumerate}

\chaptersection{Complete Problem Checklist}

The chapter's roster and the exact tool each problem needs; use this as a
one-row-per-problem review before the mixed set.

\begin{center}
\renewcommand{\arraystretch}{1.2}\setlength{\tabcolsep}{3pt}
\begin{tabular}{@{}p{0.16\linewidth} >{\raggedright\arraybackslash}p{0.24\linewidth} >{\raggedright\arraybackslash}p{0.20\linewidth} >{\raggedright\arraybackslash}p{0.10\linewidth} >{\raggedright\arraybackslash}p{0.18\linewidth}@{}}
\toprule
\textbf{ID} & \textbf{Title} & \textbf{Tool} & \textbf{Rating} & \textbf{Section}\\
\midrule
CF 2190D & Prufer Vertex & forest-completion + DP fold & 2500 & 30.2\\
CF 280C  & Game on Tree  & linearity + depth reciprocals & 2200 & 30.4\\
CF 917D  & Stranger Trees & completion + binomial inversion & 2500 & Summit\\
\bottomrule
\end{tabular}
\end{center}

\chaptersection{Recommended Study Order}

Read \(\secref{c32-prufer}\) first and hand-encode the running tree; the
bijection is the one idea the whole chapter rests on.  Then \(\secref
{c32-cayley}\), whose Cayley corollary should feel immediate.  Work the
CF 2190D sample cases by hand before reading the solution -- the
smallest-leaf peel on \(n=3\) and \(n=5\) is quick and pins the notion of
\(P(T)\) far better than a paragraph of prose.

\(\secref{c32-matrix}\) is best approached once \(\secref{c32-cayley}\) has
been digested -- the \(K_n\) sanity check on Cayley is the payoff.  Then
\(\secref{c32-expect}\), which is a short section and a fun change of
register: the same input object (a tree) but the answer is a real number.
Finish with \(\secref{c32-boundary}\), then the Summit.

\chaptersection{Difficulty Progression}

\begin{itemize}[nosep]
  \item \lv{3} Encode/decode the running tree \(T\); hand-compute \(n^{n-2}\)
        for \(n=1..5\) and verify Cayley on the shapes you can draw.
  \item \lv{4} Solve CF 280C in one DFS once you can articulate why
        \(\Pr[A_v]=1/\mathrm{depth}(v)\).  Read Matrix-Tree on \(C_5\),
        verify \(\tau(C_5)=5\) both combinatorially and by hand.
  \item \lv{5} Reconstruct the CF 2190D samples by hand and code the
        forest-completion formula.  Read the Summit's DP structure, and
        verify \(g(0)=0\) on the star \(n=3\) using the inversion.
  \item \lv{6} Solve the Summit CF 917D on your own; hand-code the
        \((k,s)\) DP up to \(n\le 20\) and check \(f(k)\) matches the
        binomial-inversion output.
\end{itemize}

\chaptersection{Reduction Gallery}

\begin{itemize}[nosep]
  \item \emph{Count spanning trees of \(K_n\)} \(\to\) Cayley \(n^{n-2}\)
        \(\to\) Pr\"ufer: sequences of length \(n-2\) over \([n]\).
  \item \emph{Count completions of a forest} \(\to\) contract each
        component to a super-vertex, then Pr\"ufer on the super-tree with
        a local \(s_i\)-way graft per super-vertex.
  \item \emph{Count labelled trees sharing exactly \(k\) edges with a
        given tree \(T\)} \(\to\) inclusion-exclusion / binomial inversion
        on top of "count containing a fixed subset of \(k\) edges"
        \(\to\) forest-completion per subset.
  \item \emph{Expected steps to empty a rooted tree by random subtree
        deletion} \(\to\) sum of indicators \(\mathbf{1}_{A_v}\)
        \(\to\) uniform-permutation symmetry on the root-\(v\) path
        \(\to\) \(1/\mathrm{depth}(v)\) per vertex.
  \item \emph{Count spanning trees of a specific graph \(G\)} \(\to\)
        Laplacian \(L(G)\) \(\to\) delete row/col \(i\) \(\to\)
        modular Gaussian elimination \(\to\) determinant.
\end{itemize}

\chaptersection{Mixed Problem Set}

\begin{mixedset}
\practice{1605D}{Treelabeling}{CF 1605D}{https://codeforces.com/problemset/problem/1605/D}{CF 2100}{greedy, tree DP, label assignment}{Assign letters to tree vertices to minimise the sum along paths; the trick is not Prüfer but a careful greedy on centroid-like structure. Read as a warm-up on tree-DP counting before the harder items.}
\practice{1311E}{Construct the Binary Tree}{CF 1311E}{https://codeforces.com/problemset/problem/1311/E}{CF 2200}{construction, tree DP, level assignment}{Construct a labelled binary tree from a level array; the count is a multinomial over level sizes once you fix the parent-choice per node -- same shape of argument as Prüfer-degree constraints.}
\practice{1540B}{Tree Array}{CF 1540B}{https://codeforces.com/problemset/problem/1540/B}{CF 2300}{tree DP, segment tree, combinatorics}{A tree-DP fold with a segment tree on the Euler tour; cross-cite Ch.19 for the monoid.}
\practice{1929F}{Sasha and the Wedding BST}{CF 1929F}{https://codeforces.com/problemset/problem/1929/F}{CF 2300}{BST count, Catalan-like, tree enumeration}{Count BST labelings on a fixed structure; the multinomial and Catalan-like recurrence are cousins of Cayley's argument.}
\practice{653E}{Bear and Forgotten Tree 2}{CF 653E}{https://codeforces.com/problemset/problem/653/E}{CF 2400}{spanning trees, Matrix-Tree, mod arithmetic}{Direct application of \(\secref{c32-matrix}\): count spanning trees of a given graph modulo a prime via Laplacian cofactor.}
\practice{1237E}{Balanced Binary Search Trees}{CF 1237E}{https://codeforces.com/problemset/problem/1237/E}{CF 2400}{tree counting, DP on sizes}{Count balanced BSTs; recursive size-DP mirrors Cayley's derivation but with a height-balanced filter.}
\practice{1120D}{Power Tree}{CF 1120D}{https://codeforces.com/problemset/problem/1120/D}{CF 2500}{Prüfer-style, MST-style greedy, tree reconstruction}{Given a distance matrix, reconstruct the tree with minimum weight; use Prüfer-like degree reasoning plus a Kruskal-style filter. Cite Ch.31.}
\practice{804D}{Expected diameter of a tree}{CF 804D}{https://codeforces.com/problemset/problem/804/D}{CF 2500}{linearity of expectation, edge indicators}{Split diameter into edge indicators; the technique is exactly \(\secref{c32-expect}\) but with edge indicators on path-uses.}
\practice{1762E}{Tree Sum}{CF 1762E}{https://codeforces.com/problemset/problem/1762/E}{CF 2600}{subtree reweight DP, tree DP}{Reweight subtree contributions and recompute the maximum path sum; a fold-heavy Ch.13 style problem with an enumerative flavour.}
\practice{1495D}{BFS Trees}{CF 1495D}{https://codeforces.com/problemset/problem/1495/D}{CF 2600}{BFS tree counting, combinatorics}{Count BFS-tree realizations of a given graph level structure; the per-level product mirrors the forest-completion \(\prod s_i\) factor.}
\end{mixedset}

\chaptersection{The Summit -- CF 917D Stranger Trees}

\begin{notebox}[Why this is the chapter top rope]
CF 917D asks a counting question whose surface is unlabelled ("share
exactly \(k\) edges with a given tree") but whose answer space is
\emph{labelled trees on \([n]\)}.  The tool is \(\thmref{c32-forest}\) (the
forest-completion formula from \(\secref{c32-cayley}\)) applied per subset
of the given tree's edges, combined with a \((k,s)\)\emph{-shaped} tree DP
that folds the product \(\prod_i s_i\) across subsets, plus a binomial
inversion to move from "at least" to "exactly".  Three tools from this
chapter (the bijection's completion corollary, a tree DP, and modular
arithmetic) and one from Ch.~13 (the bottom-up fold \(\chref{c10}\)) make
the full pipeline.
\end{notebox}

\solvedproblem{c32-summit}{Stranger Trees}{CF 917D}{https://codeforces.com/problemset/problem/917/D}{CF 2500 L5}{Pruefer bijection, forest-completion count \(\thmref{c32-forest}\), tree DP over subsets, binomial inversion}

\begin{problemstatement}[CF 917D, verified against the live CF problem page]
A labelled tree \(T\) is given on \(n\) vertices with \(1\le n\le 100\); the
next \(n-1\) lines describe its edges as pairs \((v,u)\) with \(1\le v,u\le
n\), \(v\ne u\).  For each \(k\in[0,n-1]\), count the number of
\emph{labelled} trees \(T'\in\mathcal{T}_n\) that share \emph{exactly} \(k\)
edges with \(T\) -- two labelled trees are considered different iff their
edge sets differ.  Print the \(n\) answers modulo \(10^9+7\).  Time limit
1 second, memory 256 MB.  Samples: (i) star \(1\!-\!2,\, 1\!-\!3\) on
\(n=3\) prints \(0\;2\;1\); (ii) path \(1\!-\!2,\, 2\!-\!3,\, 3\!-\!4\) on
\(n=4\) prints \(1\;7\;7\;1\); (iii) star \(1\!-\!2,\, 1\!-\!3,\, 1\!-\!4\)
on \(n=4\) prints \(0\;9\;6\;1\).
\end{problemstatement}

\begin{figure}[H]
\centering
\begin{tikzpicture}[diagram, scale=0.9, transform shape,
   tnode/.style={circle, draw=primaryblue!70, fill=primaryblue!8, line width=0.7pt,
                 minimum size=6.5mm, font=\scriptsize\bfseries, inner sep=1pt}]
  \node[tnode] (c) at (0,0) {1};
  \node[tnode] (l1) at (-1.4,1.2) {2};
  \node[tnode] (l2) at (1.4,1.2) {3};
  \draw[primaryblue!70, line width=1pt] (c)--(l1) (c)--(l2);
  \node[annot, font=\scriptsize] at (0,-1.2) {star on 3: answers 0 2 1};
  \node[tnode] (p1) at (4,1.2) {1};
  \node[tnode] (p2) at (5,0.4) {2};
  \node[tnode] (p3) at (6,-0.4) {3};
  \node[tnode] (p4) at (7,-1.2) {4};
  \draw[primaryblue!70, line width=1pt] (p1)--(p2)--(p3)--(p4);
  \node[annot, font=\scriptsize] at (5.5,-2.1) {path on 4: answers 1 7 7 1};
\end{tikzpicture}
\caption{Two of the three official samples; the third is the four-leaf star,
whose histogram is \(0\;9\;6\;1\).  The Summit recovers all three via the Solution step's pipeline.}
\label{fig:c32-summit}
\end{figure}

\begin{attemptgate}[Attempt it yourself first -- \lv{5}, give yourself 60 minutes]
Before coding: what quantity of the forest \((V,S)\) for a subset \(S\subseteq
E(T)\) controls the number of completions that contain \(S\)?  And how do
you convert "count completions containing \(S\)" into "count completions
sharing \emph{exactly} \(k\) edges with \(T\)"?
\peek{\(\thmref{c32-forest}\) gives the count of trees containing a fixed
forest \((V,S)\) directly -- the answer depends on the sizes of \((V,S)\)'s
components, not their shape.  For the second question: binomial inversion
between "at least" and "exactly" counts.}
\nudge{Define \(f(k)=\sum_{S\subseteq E(T),|S|=k}\) (trees containing \(S\)).
Define \(g(j)=\) the number of
trees sharing exactly \(j\) edges with \(T\).  Then
\(f(k)=\sum_{j\ge k}\binom{j}{k}g(j)\), because each tree sharing \(j\) edges
contributes \(\binom{j}{k}\) subsets of size \(k\).  Invert.
}
\fullhint{For \(f(k)\): iterate over \(S\subseteq E(T)\) of size \(k\),
which is a forest with \(n-k\) components of sizes \(s_1,\dots,s_{n-k}\);
\(\thmref{c32-forest}\) yields \(n^{n-k-2}\prod_i s_i\) completions per
subset.  Aggregate over all \(S\) of size \(k\) with a tree DP: root \(T\)
at \(1\), and let \(dp[u][j][t]\) be the sum over size-\(j\) subsets of
edges inside \(u\)'s subtree of the product of closed component sizes,
with the component containing \(u\) currently having size \(t\) (still
open).  Merge child subtrees by either dropping the connecting edge (the
child's open becomes closed, multiply by its open size) or keeping it
(the opens merge, \(t\) adds).  Final answer: \(f(k)=n^{n-k-2}\sum_t t\cdot
dp[1][k][t]\).  Then invert binomially.}
\end{attemptgate}

\solstep{Understanding}
The Summit is a binomial inversion wearing a Pr"ufer hat.  The natural
quantity to enumerate is \(f(k)\): the total over all \(k\)-subsets \(S\subseteq
E(T)\) of the number of labelled trees that \emph{contain} \(S\).  Each
such tree \(T'\) sharing exactly \(j\) edges with \(T\) contributes to
\(f(k)\) once per \(k\)-subset of its shared \(j\)-edge intersection, i.e.\
\(\binom{j}{k}\) times.  Summing gives \(f(k)=\sum_{j\ge k}\binom{j}{k}
g(j)\), and the standard binomial inversion recovers \(g(k)=\sum_{j\ge k
}(-1)^{j-k}\binom{j}{k}f(j)\).  The problem reduces to computing \(f(k)\)
for \(k=0,\dots,n-1\), which \(\thmref{c32-forest}\) makes a subtree-DP
exercise on \(T\) itself.

\solstep{Key Observations}
\begin{observation}
\(f(k)\) depends on \(S\) only through the \emph{component sizes} of the
forest \((V,S)\): the completion count \(n^{n-k-2}\prod_i s_i\) is
invariant under tree-shape changes that preserve the size multiset.
\end{observation}
\begin{observation}
The DP \(dp[u][j][t]\) has \(O(n^3)\) states per node (bounded by
subtree size), and the merge over a child's DP is \(O(n^2)\), so total
time is \(O(n^3)\), which fits \(n\le 100\) with room.
\end{observation}
\begin{observation}
Binomial inversion uses only \(\binom{n}{k}\bmod p\) with \(p=10^9+7\); precompute
factorials and inverse factorials in \(O(n)\) and query in \(O(1)\).
\end{observation}

\solstep{Think It With Me}
\begin{thinkbox}
\deadend{Enumerate all \(S\subseteq E(T)\)}{\(2^{n-1}\) subsets is \(2^{99}\)
for \(n=100\), hopeless.}
\deadend{Compute \(g(k)\) directly via a filter DP}{The "exactly \(k\)\
" filter requires tracking which tree edges are shared with \(T'\), but
\(T'\) is not explicit -- only its overlap with \(T\) matters, and there's
no way to enumerate \(T'\)'s edges without leaving labelled-tree space.}
\deadend{Use Matrix-Tree on a modified graph}{Matrix-Tree counts spanning
trees of a fixed \(G\), which does not match the question "count \(T'\)\
sharing exactly \(k\) edges with \(T\)" -- you would need to iterate over
exponentially many \(G\)'s.}
\aha{\(f(k)\) is a tree DP on \(T\) itself, with state \((j,t)\) per node
and a merge that either closes the child's open component or fuses it into
the parent's open.  Then binomial inversion converts \(f\) to \(g\).}
\end{thinkbox}

\solstep{The Approach Ladder}
\begin{center}\renewcommand{\arraystretch}{1.25}\setlength{\tabcolsep}{3pt}
\begin{tabular}{@{}c p{0.26\linewidth} p{0.30\linewidth} l c@{}}
\toprule
\textbf{\#} & \textbf{Approach} & \textbf{Idea} & \textbf{Time} & \textbf{Verdict}\\
\midrule
1 & \lv{2}\ enumerate all \(T'\) & \(n^{n-2}\) candidates & \(\bigO{n^{n}}\) & \tle\\
2 & \lv{3}\ enumerate all subsets \(S\subseteq E(T)\) & \(2^{n-1}\) subsets & \(\bigO{2^n}\) & \tle\\
3 & \lv{4}\ Matrix-Tree on every subgraph & exponential \(G\)'s & \(\bigO{2^n n^3}\) & \tle\\
4 & \lv{5}\ \textbf{tree DP + binomial inversion} & \((j,t)\) state per node & \(\bigO{n^3}\) & \ac\\
\bottomrule
\end{tabular}\end{center}
Rungs 1--3 all blow up on the exponential state space.  Rung 4 is a
tree-DP fold and an \(O(n^2)\) inversion, both \(\bigO{n^3}\) end to end.

\solstep{Solution}
Root \(T\) at vertex \(1\).  Compute the DP \(dp[u][j][t]\) as described.
\(f(k)=n^{n-k-2}\sum_t t\cdot dp[1][k][t]\bmod p\).  Precompute \(\binom{n}{k}\)
and \((-1)^k\) tables.  Invert binomially: \(g(k)=\sum_{j\ge k}(-1)^{j-k
}\binom{j}{k}f(j)\).  Print \(g(0),\dots,g(n-1)\).

\solstep{Implementation}
\begin{lstlisting}[caption={CF 917D. Tree DP with state (edges kept, open size) then binomial inversion. Handles the k=n-1 case by special-casing the negative exponent via modPow with MOD minus 2. ASCII only, identifiers without underscores}, label={lst:c32-summit}]
#include <bits/stdc++.h>
using namespace std;
typedef long long ll;
const ll MOD=1000000007;
ll modPow(ll b, ll e){ ll r=1; b%=MOD; while(e>0){ if(e&1) r=r*b%MOD; b=b*b%MOD; e>>=1;} return r; }
int n;
vector<vector<int> > adj;
vector<vector<vector<ll> > > dp; // dp[u][j][t]
vector<int> subSz;
void dfs(int u, int p){
  subSz[u]=1;
  dp[u][0][1]=1;
  int used=0;
  for(int v : adj[u]) if(v!=p){
    dfs(v,u);
    vector<vector<vector<ll> > > nDp(n+1, vector<vector<ll> >(n+1, vector<ll>(n+1,0)));
    for(int j1=0; j1<=used; j1++)
      for(int t1=1; t1<=subSz[u]; t1++) if(dp[u][j1][t1])
        for(int j2=0; j2<=subSz[v]-1; j2++)
          for(int t2=1; t2<=subSz[v]; t2++) if(dp[v][j2][t2]){
            ll ways=dp[u][j1][t1]*dp[v][j2][t2]%MOD;
            // drop edge (u,v): v's open component becomes closed -> * t2
            nDp[u][j1+j2][t1] = (nDp[u][j1+j2][t1] + ways*t2)%MOD;
            // keep edge (u,v): opens merge
            nDp[u][j1+j2+1][t1+t2] = (nDp[u][j1+j2+1][t1+t2] + ways)%MOD;
          }
    for(int j=0;j<=n;j++) for(int t=0;t<=n;t++) dp[u][j][t]=nDp[u][j][t];
    used += subSz[v];
    subSz[u]+=subSz[v];
  }
}
int main(){
  if(scanf("%d",&n)!=1) return 0;
  adj.assign(n+1, {});
  for(int i=0;i<n-1;i++){ int u,v; scanf("%d %d",&u,&v); adj[u].push_back(v); adj[v].push_back(u); }
  dp.assign(n+1, vector<vector<ll> >(n+1, vector<ll>(n+1,0)));
  subSz.assign(n+1,0);
  dfs(1,0);
  vector<ll> f(n+1,0);
  for(int k=0;k<=n-1;k++){
    ll sum=0;
    for(int t=1;t<=n;t++) sum=(sum + (ll)t*dp[1][k][t])%MOD;
    ll c=(k==n-1)?1:modPow(n, n-k-2);
    f[k]=sum*c%MOD;
  }
  // binomial inversion g(k) = sum_{j>=k} (-1)^{j-k} C(j,k) f(j)
  vector<ll> fact(n+1,1), invFact(n+1,1);
  for(int i=1;i<=n;i++) fact[i]=fact[i-1]*i%MOD;
  invFact[n]=modPow(fact[n], MOD-2);
  for(int i=n;i>=1;i--) invFact[i-1]=invFact[i]*i%MOD;
  auto nCk=[&](int a,int b)->ll{ if(b<0||b>a) return 0; return fact[a]*invFact[b]%MOD*invFact[a-b]%MOD; };
  for(int k=0;k<=n-1;k++){
    ll gk=0;
    for(int j=k;j<=n-1;j++){
      ll term=nCk(j,k)*f[j]%MOD;
      if((j-k)&1) gk=(gk-term+MOD)%MOD; else gk=(gk+term)%MOD;
    }
    printf("%lld ", gk);
  }
  printf("\n");
  return 0;
}
\end{lstlisting}

\solstep{Trace on the Official Samples}
\begin{tracebox}[Sample 1: \(n=3\) star \(1\!-\!2,\,1\!-\!3\)]
DP gives \(f(0)=3^{1}\cdot 1\cdot 1\cdot 1=3\) (empty subset; three isolates;
\(\prod s_i=1\)).  \(f(1)\) sums over two 1-edge subsets, each giving
components \((2,1)\) with \(\prod s_i=2\), times \(3^{0}=1\): \(f(1)=2+2=4\).
\(f(2)=1\) (the whole star; \(c=1\); the exponent \(-1\) is special-cased
with \(\prod s_i=3\), giving \(3\cdot 3^{-1}=1\)).  Inversion: \(g(0)=f(0)-f(1
)+f(2)=3-4+1=0\); \(g(1)=f(1)-2f(2)=4-2=2\); \(g(2)=f(2)=1\).  Output
\(0\;2\;1\), matching the sample.
\end{tracebox}
\begin{tracebox}[Sample 2: \(n=4\) path \(1\!-\!2,\,2\!-\!3,\,3\!-\!4\)]
\(f(0)=4^{2}=16\); \(f(1)=3\cdot 8=24\) (three 1-edge subsets, each \(c=3\), \(\prod s_i=2\), times \(4^{1}=4\), giving \(8\) per subset); \(f(2)=
3+3+4=10\) (two adjacent-edge subsets give \(c=2,\prod s_i=3\) each; one
disjoint subset gives \(\prod s_i=4\)); \(f(3)=1\).  Inversion: \(g(0
)=16-24+10-1=1\); \(g(1)=24-20+3=7\); \(g(2)=10-3=7\); \(g(3)=1\).  Output \(1\;7\;7\;1\),
matching the sample.
\end{tracebox}
\begin{tracebox}[Sample 3: \(n=4\) star \(1\!-\!2,\,1\!-\!3,\,1\!-\!4\)]
\(f(0)=16\); \(f(1)=24\) (three 1-edge subsets, same as Sample 2); \(f(2
)=9\) (all three 2-edge subsets are the two-leaf star, \(c=2\) and \(\prod s_i=3\) each); \(f(3)=1\).  Inversion: \(g(0)=16-24+9-1=0\); \(g(1)=24-18+3=9\); \(g(2)=9-3=6\); \(g(3)=1\).  Output
\(0\;9\;6\;1\), matching the sample.
\end{tracebox}

\solstep{Complexity}
\begin{complexitybox}
The DFS visits \(O(n)\) nodes; the child merge at node \(u\) is \(O((s_u)\cdot s_v)\) states on each side, so the total across the tree is
\(O(n^3)\).  Binomial inversion is \(O(n^2)\).  Space \(O(n^3)\) for the
DP array, \(O(n)\) for the precomputed factorials.  For \(n=100\) the
DP is \(\sim 10^6\) states, comfortably inside 1 second.  All arithmetic
mod \(10^9+7\); use 64-bit throughout.
\end{complexitybox}

\begin{hardeningbox}
\begin{itemize}[nosep,leftmargin=1.4em]
\item \textbf{The \(k=n-1\) exponent.}  \(n^{n-k-2}\) reads as \(n^{-1}\)
when \(k=n-1\); special-case to \(1\) (already a tree, one completion).
\item \textbf{Reset the DP.}  \(dp[u][j][t]\) is per-node; allocate
\((n+1)^3\) once and initialise to zero, or use a per-node vector.
\item \textbf{Child loop bound.}  \(j_2\) runs over \([0, s_v-1]\), not
\([0,n]\); tightening this bound is the difference between \(\bigO{n^4}\)
and \(\bigO{n^3}\).
\item \textbf{Sign on the inversion.}  \((-1)^{j-k}\) is a mod-subtraction;
always write \((x-y+\text{MOD})\%\text{MOD}\) to keep the value nonnegative.
\item \textbf{Rooted at 1.}  The DP assumes a root; picking a different
root permutes which open size at \(u\) corresponds to which component,
but the final sum over \(t\cdot dp[1][k][t]\) is invariant -- still,
fix a root for the code.
\end{itemize}
\end{hardeningbox}

\begin{recallbox}[Chapter 30 in ten seconds]
\emph{Pr"ufer is the bijection; Cayley is the corollary; the forest-
completion formula is the workhorse; Matrix-Tree hands you a determinant
when the graph is fixed; linearity of expectation turns an exponential
state process into a sum of reciprocals of depths.  CF 2190D reads
Pr"ufer-style per vertex, CF 280C is one DFS on depth reciprocals, and
the Summit CF 917D is a tree DP with a binomial inversion.}  The
chapter ends with the same object it started with -- a labelled tree on
\([n]\) -- but you can now \emph{count} them under any reasonable filter
and \emph{expect} things about them without writing a recursion.
\selftest{
  \item State the completion formula \(\thmref{c32-forest}\) and its
        \(k=n\) and \(k=1\) special cases.
  \item Give the binomial-inversion identity relating \(f(k)\) and \(g(k)\).
  \item Explain why a naive \(P(T)=c_{n-2}\) filter fails on CF 2190D.
  \item Reconstruct the running example \(n=3\) star Summit pipeline
        from memory.
}
\end{recallbox}

\chaptersection{Tool Kit Inventory}

\begin{center}
\renewcommand{\arraystretch}{1.2}\setlength{\tabcolsep}{3pt}
\begin{tabular}{@{}p{0.28\linewidth} >{\raggedright\arraybackslash}p{0.22\linewidth} >{\raggedright\arraybackslash}p{0.36\linewidth}@{}}
\toprule
\textbf{Tool} & \textbf{Owned?} & \textbf{Use in this chapter}\\
\midrule
Pr"ufer encode/decode & Yes, \(\secref{c32-prufer}\) & the whole chapter's frame \\
Cayley \(n^{n-2}\) & Yes, \(\thmref{c32-cayley}\) & base count \\
Rooted \(n^{n-1}\) & Yes, \(\thmref{c32-rooted}\) & normaliser \\
Degree multinomial & Yes, \(\corref{c32-mult}\) & CF 2190D filter \\
Forest completion & Yes, \(\thmref{c32-forest}\) & CF 2190D, Summit \\
Laplacian + Matrix-Tree & Yes, \(\secref{c32-matrix}\) & given-graph counts \\
Modular determinant & Yes, \(\coderef{c32-moddet}\) & Matrix-Tree engine \\
Linearity of expectation & Yes, \(\lemref{c32-indexp}\) & CF 280C \\
Binomial inversion & Yes, Summit & \(f\to g\) filter \\
Depth from root & no, Ch.12 (\chref{c09}) & \(\secref{c32-expect}\) reads it \\
Tree DP fold & no, Ch.13 (\chref{c10}) & Summit merges via it \\
DSU component census & no, Ch.17 (\chref{c28}) & CF 2190D forest read \\
Segment tree monoid & no, Ch.19 (\chref{c13}) & cited once, not required here \\
AHU tree isomorphism & no, Ch.28 (\chref{c22}) & \(\secref{c32-boundary}\) hand-off \\
\bottomrule
\end{tabular}
\end{center}

\chaptersection{Related Topics}

\begin{itemize}[nosep]
  \item Ch.~13 (\(\chref{c10}\)) tree DP -- the bottom-up fold; the Summit's
        \((j,t)\) DP is that fold with an "open component size" state
        layered on top.
  \item Ch.~17 (\(\chref{c28}\)) DSU and small-to-large -- CF 2190D reads
        the forest's components via DSU; the fold's merge can also be
        sped up via small-to-large.
  \item Ch.~26 (\(\chref{c30}\)) virtual trees -- the handoff in
        \(\secref{c32-boundary}\) when the query is over a marked vertex
        subset rather than the whole \([n]\).
  \item Ch.~28 (\(\chref{c22}\)) tree isomorphism and AHU -- the unlabelled
        side of \(\thmref{c32-orbit}\); needed to compute \(|\mathrm{aut}(S)|\) when reconciling labelled and unlabelled counts.
  \item Ch.~29 (\(\chref{c31}\)) trees from graphs -- the previous chapter;
        it ends where Ch.~30 starts (a tree is given); Ch.~30 begins at
        the count.
  \item Ch.~31 (\(\chref{c17}\)) MST and variants -- the next chapter;
        a different "tree inside a graph" that counts via a greedy, not
        a determinant.
  \item Ch.~12 (\(\chref{c09}\)) tree metrics -- depth array used verbatim
        by \(\secref{c32-expect}\).
  \item Ch.~19 (\(\chref{c13}\)) segment trees -- cited for the range-max
        monoid used in Chapter 29's Summit; no segment tree is needed
        inside Chapter 30 itself.
\end{itemize}

\begin{lastword}[Last Word --- Chapter 30]
Chapter~30 closes the enumerative loop opened in Chapter~1: the same
labelled tree you first met as an input object is here a \emph{counting
object}, and you own four ways to answer questions about the family.  The
Pr"ufer bijection is the deepest; every other tool in this chapter is
either a corollary of it (Cayley, the completion formula) or a sibling
engine for the case it does not cover (Matrix-Tree for a given graph,
linearity for an expectation).  The Summit CF 917D is a deliberate
capstone: it uses three of the four tools in one pipeline, and the fourth
(matrix-tree) is what you would reach for if the input were not a \emph{
tree} but a graph.  The next chapter, Ch.~31 (\(\chref{c17}\)), turns the
"given graph" branch into a weighted one and hands you Kruskal, Prim, and
the reconstruction tree -- the optimisation twin of the enumeration you
have just finished.
\end{lastword}

